\documentclass[10pt,a4paper]{amsart}
\usepackage[margin=19mm]{geometry}
\usepackage{amsmath,amssymb,mathtools}
\usepackage[scr=rsfs,cal=euler]{mathalfa}
\usepackage{xfrac}
\usepackage[unicode,hidelinks,bookmarks=false]{hyperref}
\newcommand{\ie}{i.e.\ }
\newcommand{\cf}{cf.\ }
\newcommand{\resp}{respectively\ }
\newcommand{\Subsection}{\subsection}
\providecommand{\nobreakdash}{\nobreak\hbox{-}}
\setlength{\emergencystretch}{2em}
\multlinegap0pt
\usepackage{amssymb}
\usepackage{tikz}
\usepackage{enumerate}
\newtheorem{definition}{Definition}
\newtheorem{proposition}{Proposition}
\newcommand{\Hom}{\mathrm{Hom}}
\newcommand{\Ob}{\mathrm{Ob}}
\newcommand{\Z}{\mathbb{Z}}
\newcommand{\calG}{\mathcal{G}}
\newcommand{\m}{\mathfrak{m}}     
\title{\bf Multivariate Quandles as Groupoid Invariants}
\author{Xerxes D. Arsiwalla, Louis H. Kauffman}
\date{Source: arXiv:2609.30262v1 (CC BY 4.0)}
\begin{document}
\maketitle

\noindent\emph{Source and licence.} \bf Multivariate Quandles as Groupoid Invariants. Version arXiv:2609.30262v1 (CC BY 4.0), distributed by arXiv under the Creative Commons Attribution 4.0 International licence, http://creativecommons.org/licenses/by/4.0/, as recorded in the arXiv licence metadata for this exact version and shown on the abstract page https://arxiv.org/abs/2609.30262v1. Full licence notice: \texttt{licenses/arxiv-2609.30262-CC-BY-4.0.txt}.\par\smallskip

\noindent\textit{Editorial scope.} Counted unit: the article's proposition prop-id (source0.tex lines 707--717). The article's complete original proof of it is delivered with it (source0.tex lines 718--738, one \texttt{proof} environment of 21 lines). No mathematics is written, repaired or re-derived: the statement and the proof are the article's own lines, in the article's own notation, and no retained line is substituted or re-flowed.  Retained uncounted as the source-local context the proof needs: lines 403--435.  Nothing was omitted from any retained span, so the package carries no omission record.  No retained line contains a citation, so the wrapper carries no bibliography and no bibliography entry.  No retained line contains a figure environment, an image-inclusion command or drawing code, so no image is delivered and the package declares no asset dependency.  Editorial formatting only, in addition: an AMS \texttt{amsart} preamble that restates the article's own declarations and macros used by the retained lines, namely the environments \texttt{definition}, \texttt{proposition} on separate counters, together with the standard packages those lines need.  The retained environments print in this document as Definition 1, Proposition 1 in order of appearance, whereas the article prints the same items with the numbering of its own full document.  The complete native source of the article as distributed in the arXiv e-print for this exact version is bundled unchanged as \texttt{sources/211737/source0.tex}.
\begin{definition}
\label{defgcygr}
Let $n_1, \ldots, n_m$ be positive integers. We define the \emph{groupoid}  $\calG(n_1, \ldots, n_m)$ as a category, whose objects are:  
\[ \Ob(\calG) = \{1, 2, \ldots, m\}  \]
with morphisms:  
 \[
          \Hom(i,j) =
          \begin{cases}
            \Z_{n_i} & \text{if } i = j, \\
            \varnothing & \text{if } i \neq j.
          \end{cases}
        \]
  For $a, b \in \Hom(i,i) = \Z_{n_i}$,  compositions are defined by
        \[
          b \circ a = a + b \pmod{n_i}.
        \]
 subject to associativity   
        \[
          (c \circ b) \circ a = c \circ (b \circ a) \pmod{n_i},
        \]
     for any $a, b, c \in \Z_{n_i}$.  
     
Furthermore, for each object $i$, there exists an identity morphism  $\mathrm{id}_i = 0 \in \Z_{n_i}$, such that  for any $a \in \Z_{n_i}$,
        \[
          a \circ \mathrm{id}_i =  a =  \mathrm{id}_i \circ a 
        \]

And,  for each $a \in \Hom(i,i) = \Z_{n_i}$, the inverse morphism is defined as 
        $a^{-1} = -a \pmod{n_i}$,  satisfying 
        \[
          a \circ a^{-1} = \mathrm{id}_i = a^{-1} \circ a
        \]    
\end{definition}

\begin{proposition}    
\label{prop-id}  
For  $\calG = \calG(n_1, \ldots, n_m)$ (Definition~\ref{defgcygr}), and  $M$  a $\Z[\calG]$-module, with   $M_i := e_i \cdot M$ for each  $i$, we have that:  
\begin{enumerate} 
  \item Each $M_i$ realizes a $\Z[\Z_{n_i}]$-module.
  \item The idempotents $\{ e_i \}$ uniquely specify the module decomposition  
        \[
          M \cong \bigoplus_{i=1}^{m} M_i.
        \]         
\end{enumerate}
\end{proposition}

\begin{proof}
(1) The subset $e_i M = \{e_i \cdot \m: \m\in M\}$ is an abelian
subgroup of $M$. For $\alpha \in \Z[\Z_{n_i}]$ and $e_i \m\in M_i$,
define
\[
  \alpha \cdot (e_i m) := \tilde{\alpha} \cdot \m 
\]
where $\tilde{\alpha} = (0, \ldots, \alpha, \ldots, 0) \in
\prod_j \Z[\Z_{n_j}] \cong \Z[\calG]$ with $\alpha$ in position $i$.
Since $e_i$ acts as the identity on $M_i$ and annihilates $M_j$ for
$j \neq i$, this is well-defined and gives a $\Z[\Z_{n_i}]$-module
structure.

(2) By completeness, every $\m \in M$ can be written as
\[
  \m = 1 \cdot \m= \left(\sum_{i=1}^{m} e_i\right) \cdot \m
  = \sum_{i=1}^{m} e_i \m
\]
with $e_i \m \in M_i$. Furthermore, by orthogonality, if $\sum_i \m_i = 0$ with
$\m_i \in M_i$, then applying $e_j$ gives $\m_j = 0$. Thus, $\sum_i \m_i$ is linearly independent for every $\m$, which uniquely specifies the   decomposition of $M$ as a $\Z[\calG]$-module.    
\end{proof}

\end{document}
